\chapter{The Engine Interface}

A wick program talks to the lantern engine through the `lt` namespace. Every
`lt.*` call is registered with a typed signature that the compiler checks:
wrong arity or a wrong argument type is a compile error naming the call and
the argument position, exactly as for your own functions. This chapter
describes the frame contract that structures every game, documents the API by
area, and then covers the two properties --- determinism and verifiability ---
that the engine and language were co-designed to provide.

\section{The frame contract}

The engine drives the game by calling two functions it looks up by name, once
each per frame, at 60 frames per second:

\begin{lstlisting}[language=wick]
fn update(dt: num) { }   // simulation; dt is seconds since the last frame
fn draw() { }            // rendering; 3D calls first, 2D composites on top
\end{lstlisting}

\noindent
`update(dt)` advances the simulation. Its argument `dt` is the elapsed time in
seconds, \emph{capped at 0.1} so that a long stall (a breakpoint, a slow load)
cannot produce a single enormous step that tunnels the player through a wall.
`draw()` renders the frame; 3D calls are issued first and the 2D calls
composite on top, so a HUD drawn with `lt.print` and `lt.rect` always sits
over the scene.

Both functions are optional. A `main.wick` with neither is legal --- it runs
its top-level statements once and then presents a static frame. Top-level
statements, everything not inside a `fn`, run exactly once at load, and that
is where a game creates its meshes, loads its textures and sounds, and
initialises its state. The order is fixed and worth holding in mind: top-level
once, then `update` and `draw` every frame, then collection at the frame
boundary (Chapter~7).

\section{Frame and screen}

\begin{center}
\small\begin{tabular}{@{}p{0.39\linewidth}p{0.53\linewidth}@{}}
\toprule
\textbf{Call} & \textbf{Notes} \\
\midrule
\texttt{lt.W} \enskip \texttt{lt.H} & compile-time constants: 400, 240 \\
\texttt{lt.clear(r, g, b)} & clear colour and depth; call first in \texttt{draw()} \\
\texttt{lt.time(): num} & seconds since boot (fixed-step under \texttt{LANTERN\_FIXED\_DT}) \\
\texttt{lt.screenshot(path: str)} & save the 400$\times$240 frame as BMP \\
\texttt{lt.quit()} & request a clean exit \\
\texttt{lt.escape\_quits(enable: bool)} & Escape quits by default; disable for pause menus \\
\bottomrule
\end{tabular}
\end{center}

\noindent
`lt.W` and `lt.H` are numeric constants, not calls --- they resolve at compile
time to 400 and 240 --- so `lt.W / 2` is a constant expression you can use
freely for layout. The screen is always 400 by 240; the engine scales the
presented image to the window.

\section{The 3D scene}

\begin{center}
\small\begin{tabular}{@{}p{0.43\linewidth}p{0.49\linewidth}@{}}
\toprule
\textbf{Call} & \textbf{Notes} \\
\midrule
\texttt{lt.camera(ex,ey,ez, tx,ty,tz, fov=55)} & eye position and look-at target \\
\texttt{lt.light(dx,dy,dz, ambient=0.35)} & the directional light \\
\texttt{lt.point\_light(i, x,y,z, radius, r=1,g=1,b=1)} & slots 0--3; \texttt{radius <= 0} turns the slot off \\
\texttt{lt.fog(start, end, r,g,b)} & linear by view depth; \texttt{end <= start} disables \\
\bottomrule
\end{tabular}
\end{center}

\noindent
Trailing parameters with defaults --- `fov`, `ambient` --- may be omitted at the
call site; the compiler fills the default in, so a native always receives full
arity. There are four point-light slots, indexed 0 through 3; setting a slot's
radius to zero or below disables it, which is how Kora Night extinguishes a
lamp:

\begin{lstlisting}[language=wick]
if lamp_lit[i] {
  lt.point_light(i, lamp_x[i], 0.8, lamp_z[i], 3.6, 1.0, 0.70, 0.28)
} else {
  lt.point_light(i, 0, 0, 0, 0, 0, 0, 0)   // radius 0: off
}
\end{lstlisting}

\section{Meshes}

\begin{center}
\small\begin{tabular}{@{}p{0.43\linewidth}p{0.49\linewidth}@{}}
\toprule
\textbf{Call} & \textbf{Notes} \\
\midrule
\texttt{lt.cube(): num} & unit cube \\
\texttt{lt.plane(segs=1): num} & unit XZ plane, tessellated \\
\texttt{lt.sphere(seg=12)}, \texttt{lt.cylinder(seg=16)}, \texttt{lt.cone(seg=16)} & unit primitives \\
\texttt{lt.mesh(verts: list<num>): num} & custom mesh, 12 floats per vertex \\
\texttt{lt.load\_mesh(path: str): num} & Wavefront OBJ (flat-normal fallback) \\
\texttt{lt.draw(m, x,y,z, rx,ry,rz, sx,sy,sz, r,g,b, tex=-1)} & gouraud-lit draw (rotations, scales, colour default) \\
\texttt{lt.draw\_lerp(a, b, t, x,y,z, \ldots, tex=-1)} & keyframe tween between two same-count meshes \\
\texttt{lt.billboard(tex, x,y,z, w,h, u0,v0,u1,v1)} & camera-facing quad (sprite characters) \\
\texttt{lt.shadow(x,y,z, radius, alpha=0.35)} & grounding blob \\
\bottomrule
\end{tabular}
\end{center}

\noindent
Every mesh constructor returns a `num` handle; there is no mesh type, because
handles are opaque integers the engine owns. `lt.mesh` takes a flat
`list<num>` of twelve floats per vertex --- position (3), normal (3), texture
coordinates (2), and colour (4) --- which is a natural fit for the flat-list
discipline of Chapter~5. `lt.draw` has a long tail of defaulted parameters, so
the common case stays short:

\begin{lstlisting}[language=wick]
let ground = lt.plane(12)
lt.draw(ground, 0, 0, 0, 0, 0, 0, 15, 1, 15, 0.42, 0.40, 0.40)
\end{lstlisting}

\section{2D}

\begin{center}
\small\begin{tabular}{@{}p{0.43\linewidth}p{0.49\linewidth}@{}}
\toprule
\textbf{Call} & \textbf{Notes} \\
\midrule
\texttt{lt.rect(x,y,w,h, r,g,b, a=1)} & filled, alpha-blended \\
\texttt{lt.load\_texture(path: str): num} & BMP; magenta (255,0,255) is transparent \\
\texttt{lt.sprite(tex, x,y, sx=1,sy=1)} & top-left anchored \\
\texttt{lt.sprite\_ex(tex, cx,cy, sx=1,sy=1, rot=0, r=1,g=1,b=1,a=1)} & centre anchor, rotate, tint; negative scale flips \\
\texttt{lt.sprite\_uv(tex, x,y,w,h, u0,v0,u1,v1)} & atlas sub-rect (tilemaps) \\
\texttt{lt.print(text: str, x,y, r=1,g=1,b=1,a=1)} & built-in 8$\times$8 font, \texttt{\textbackslash n} aware \\
\bottomrule
\end{tabular}
\end{center}

\noindent
`lt.print` renders in a built-in 8-by-8 pixel font and understands `\n`.
Because `+` concatenates only `str` with `str`, building a HUD line means
converting numbers explicitly with `str` --- which, unlike Lua's
`string.format`, prints integers cleanly with no trailing decimal:

\begin{lstlisting}[language=wick]
lt.print("LAMPS RELIT: " + str(score), 4, 3, 1, 1, 1, 0.95)
\end{lstlisting}

\section{Audio, input, and saves}

\begin{center}
\small\begin{tabular}{@{}p{0.43\linewidth}p{0.49\linewidth}@{}}
\toprule
\textbf{Call} & \textbf{Notes} \\
\midrule
\texttt{lt.load\_sound(path: str): num} & WAV \\
\texttt{lt.play(sound, volume=1, loop=0): num} & returns a channel, or $-1$ if all 16 are busy \\
\texttt{lt.stop(channel)} \enskip \texttt{lt.volume(v)} & \\
\texttt{lt.key(name: str): bool} & held; keyboard and gamepad merged \\
\texttt{lt.pressed(name: str): bool} & went down this frame \\
\texttt{lt.gamepad(): bool} \enskip \texttt{lt.rumble(low, high, ms)} & \\
\texttt{lt.save(name: str, data: str): bool} & binary-safe \\
\texttt{lt.load\_save(name: str): str?} & \textbf{optional} --- \texttt{nil} when unreadable \\
\bottomrule
\end{tabular}
\end{center}

\noindent
`lt.key` is true while a key is held; `lt.pressed` is true only on the frame
the key goes down --- the difference between continuous movement and a
one-shot action. The recognised input names are:

\begin{quote}
\ttfamily left\quad right\quad up\quad down\quad z\quad x\quad c\quad space\quad
return\quad escape\quad a\quad s\quad d\quad w
\end{quote}

\noindent
Keyboard and gamepad are merged onto these names: the pad's d-pad and stick map
to the directions, and its face buttons map A to `z`, B to `x`, Y to `c`, and
Start to `return`. So a game written against `lt.key("z")` works with either
input with no extra code.

The one call in this table that returns an optional is `lt.load_save`, and it
is no accident that the chapter which motivates the whole language turns on
it. Its signature is `load_save(str): str?`; the compiler will not let you use
the result until you have handled the `nil`, which is the entire point of
Chapter~4.

\section{Determinism}

wick and lantern are built so that the same program with the same inputs
produces the same frames on every machine. Two facilities make this true, and
a testing culture is built on top of them.

\paragraph{Deterministic randomness.} The built-in `rand()` returns a uniform
value in \texttt{[0, 1)} from a fixed \mbox{xorshift64$^{*}$} generator --- the
\emph{same sequence on every machine}. There is no time-seeded mode. If you
want variation between runs, you seed it yourself with `srand(seed)`, using a
number you chose and can log:

\begin{lstlisting}[language=wick]
srand(7)                    // same layout every run, on every machine
let stars: list<num> = []
for i in 0..40 {
  push(stars, rand() * 400)
  push(stars, rand() * 240)
}
\end{lstlisting}

\noindent
Because the sequence is fixed, a seeded run replays identically, which is what
makes screenshot tests meaningful --- the ``random'' starfield is the same
starfield in CI as on your desk.

\paragraph{Fixed-step time.} `lt.time()` returns seconds since boot, and under
the environment variable `LANTERN_FIXED_DT` the frame clock advances by a
fixed step regardless of wall-clock speed. Combined with deterministic
`rand()`, this makes an entire gameplay session reproducible frame for frame.

\paragraph{Screenshot capture.} Setting `LANTERN_SHOT=/tmp/x` tells the engine
to capture frame~60 as a BMP and exit. With `LANTERN_FIXED_DT=1` alongside it,
that capture is byte-identical everywhere. The two together turn ``does it
still look right'' into a mechanical check:

\begin{lstlisting}[language=bash,basicstyle=\ttfamily\small,keywordstyle={}]
LANTERN_SHOT=/tmp/x LANTERN_FIXED_DT=1 ./build/lantern games/mygame
\end{lstlisting}

\section{A culture of verification}

Determinism is only useful if the game can drive and check itself, and wick
gives two built-ins for exactly that.

\paragraph{Self-play through \texttt{env}.} `env(name): str?` reads an
environment variable --- an optional, `nil` when unset. A game can expose its
own self-play switch and branch on it, so CI runs the game on autopilot with a
fixed seed:

\begin{lstlisting}[language=wick]
let AUTO = env("SHOWCASE_AUTO") != nil   // deterministic self-play (CI)
if AUTO { srand(7) }
\end{lstlisting}

\noindent
In Kora Night, `AUTO` drives the player toward the nearest dark lamp instead
of reading the keyboard, so a headless CI run plays a real, deterministic
round and the frame-60 screenshot is a true regression test.

\paragraph{Assertions with \texttt{check}.} `check(cond: bool, msg: str)` is a
runtime assertion: when `cond` is false it stops the frame and shows
\texttt{check failed: <msg>} on the error screen. The wick test suite is built
from these --- a scene that sets up a known state and `check`s the results,
run under fixed-step time so the outcome is exact:

\begin{lstlisting}[language=wick]
check(len(lamp_x) == 4, "expected four lamps")
check(nearest_unlit() == -1, "all lamps should start lit")
\end{lstlisting}

\noindent
Together, `env`-driven self-play, deterministic `rand`, fixed-step time, and
`check` turn a game into something that tests itself the same way on every
machine --- the practical payoff of designing the language and the engine for
determinism from the start.
